\documentclass[11pt]{article}
\usepackage[a4paper,margin=0.8in]{geometry}
\usepackage{calc}
\usepackage{array,longtable,booktabs}
\usepackage{xcolor}
\usepackage{hyperref}
\usepackage{enumitem}
\usepackage{fancyvrb}
\usepackage{fvextra}
\usepackage{framed}
\usepackage{color}
\fvset{breaklines=true, breakanywhere=true}
\DefineVerbatimEnvironment{Highlighting}{Verbatim}{breaklines,breakanywhere,commandchars=\\\{\}}
\usepackage{color}
\usepackage{fancyvrb}
\newcommand{\VerbBar}{|}
\newcommand{\VERB}{\Verb[commandchars=\\\{\}]}
\DefineVerbatimEnvironment{Highlighting}{Verbatim}{commandchars=\\\{\}}
% Add ',fontsize=\small' for more characters per line
\usepackage{framed}
\definecolor{shadecolor}{RGB}{248,248,248}
\newenvironment{Shaded}{\begin{snugshade}}{\end{snugshade}}
\newcommand{\AlertTok}[1]{\textcolor[rgb]{0.94,0.16,0.16}{#1}}
\newcommand{\AnnotationTok}[1]{\textcolor[rgb]{0.56,0.35,0.01}{\textbf{\textit{#1}}}}
\newcommand{\AttributeTok}[1]{\textcolor[rgb]{0.13,0.29,0.53}{#1}}
\newcommand{\BaseNTok}[1]{\textcolor[rgb]{0.00,0.00,0.81}{#1}}
\newcommand{\BuiltInTok}[1]{#1}
\newcommand{\CharTok}[1]{\textcolor[rgb]{0.31,0.60,0.02}{#1}}
\newcommand{\CommentTok}[1]{\textcolor[rgb]{0.56,0.35,0.01}{\textit{#1}}}
\newcommand{\CommentVarTok}[1]{\textcolor[rgb]{0.56,0.35,0.01}{\textbf{\textit{#1}}}}
\newcommand{\ConstantTok}[1]{\textcolor[rgb]{0.56,0.35,0.01}{#1}}
\newcommand{\ControlFlowTok}[1]{\textcolor[rgb]{0.13,0.29,0.53}{\textbf{#1}}}
\newcommand{\DataTypeTok}[1]{\textcolor[rgb]{0.13,0.29,0.53}{#1}}
\newcommand{\DecValTok}[1]{\textcolor[rgb]{0.00,0.00,0.81}{#1}}
\newcommand{\DocumentationTok}[1]{\textcolor[rgb]{0.56,0.35,0.01}{\textbf{\textit{#1}}}}
\newcommand{\ErrorTok}[1]{\textcolor[rgb]{0.64,0.00,0.00}{\textbf{#1}}}
\newcommand{\ExtensionTok}[1]{#1}
\newcommand{\FloatTok}[1]{\textcolor[rgb]{0.00,0.00,0.81}{#1}}
\newcommand{\FunctionTok}[1]{\textcolor[rgb]{0.13,0.29,0.53}{\textbf{#1}}}
\newcommand{\ImportTok}[1]{#1}
\newcommand{\InformationTok}[1]{\textcolor[rgb]{0.56,0.35,0.01}{\textbf{\textit{#1}}}}
\newcommand{\KeywordTok}[1]{\textcolor[rgb]{0.13,0.29,0.53}{\textbf{#1}}}
\newcommand{\NormalTok}[1]{#1}
\newcommand{\OperatorTok}[1]{\textcolor[rgb]{0.81,0.36,0.00}{\textbf{#1}}}
\newcommand{\OtherTok}[1]{\textcolor[rgb]{0.56,0.35,0.01}{#1}}
\newcommand{\PreprocessorTok}[1]{\textcolor[rgb]{0.56,0.35,0.01}{\textit{#1}}}
\newcommand{\RegionMarkerTok}[1]{#1}
\newcommand{\SpecialCharTok}[1]{\textcolor[rgb]{0.81,0.36,0.00}{\textbf{#1}}}
\newcommand{\SpecialStringTok}[1]{\textcolor[rgb]{0.31,0.60,0.02}{#1}}
\newcommand{\StringTok}[1]{\textcolor[rgb]{0.31,0.60,0.02}{#1}}
\newcommand{\VariableTok}[1]{\textcolor[rgb]{0.00,0.00,0.00}{#1}}
\newcommand{\VerbatimStringTok}[1]{\textcolor[rgb]{0.31,0.60,0.02}{#1}}
\newcommand{\WarningTok}[1]{\textcolor[rgb]{0.56,0.35,0.01}{\textbf{\textit{#1}}}}
\setlist[itemize]{leftmargin=*}
\setlist[enumerate]{leftmargin=*}
\hypersetup{colorlinks=true,linkcolor=blue,urlcolor=blue}
\providecommand{\tightlist}{%
  \setlength{\itemsep}{0pt}\setlength{\parskip}{0pt}}
\title{R Workshop -- Day 1 -- Getting Comfortable with R}
\author{Applied R for Statistical Teaching, Research and Reproducible Analysis}
\date{}
\begin{document}
\maketitle
\tableofcontents
\newpage
\hypertarget{r-workshop-day-1-getting-comfortable-with-r}{%
\section{R Workshop -- Day 1 -- Getting Comfortable with
R}\label{r-workshop-day-1-getting-comfortable-with-r}}

\textbf{Session 1 of 7 \textbar{} Duration: 3 hours} \textbf{Programme:}
Applied R for Statistical Teaching, Research and Reproducible Analysis

\hypertarget{learning-objectives}{%
\subsection{Learning objectives}\label{learning-objectives}}

By the end of this session you will be able to:

\begin{enumerate}
\def\labelenumi{\arabic{enumi}.}
\tightlist
\item
  Tell R, RStudio, scripts and the console apart, and know which one to
  use for what.
\item
  Start every piece of work from an RStudio Project, using relative
  paths instead of hard-coded folder names.
\item
  Create R objects with assignment, and write comments that explain
  \emph{why}, not \emph{what}.
\item
  Create, inspect and index vectors using positions and logical
  conditions.
\item
  Represent missing values with \texttt{NA} and use \texttt{na.rm}
  correctly.
\item
  Build a small data frame and compute descriptive statistics, including
  grouped summaries.
\item
  Read an R error or warning message and find the relevant help page.
\end{enumerate}

\hypertarget{capstone-link}{%
\subsection{Capstone link}\label{capstone-link}}

Day 1 has no anchor dataset yet -- that begins on Day 2. Today you
complete a \textbf{baseline task} that you will repeat, with an
equivalent dataset, on Day 7. The comparison between your Day 1 and Day
7 work is your own evidence of progress through the workshop.

\hypertarget{segment-plan-3-hours}{%
\subsection{Segment plan (3 hours)}\label{segment-plan-3-hours}}

\begin{longtable}[]{@{}
  >{\raggedright\arraybackslash}p{(\columnwidth - 4\tabcolsep) * \real{0.3000}}
  >{\raggedleft\arraybackslash}p{(\columnwidth - 4\tabcolsep) * \real{0.4000}}
  >{\raggedright\arraybackslash}p{(\columnwidth - 4\tabcolsep) * \real{0.3000}}@{}}
\toprule\noalign{}
\begin{minipage}[b]{\linewidth}\raggedright
Segment
\end{minipage} & \begin{minipage}[b]{\linewidth}\raggedleft
Time
\end{minipage} & \begin{minipage}[b]{\linewidth}\raggedright
Purpose
\end{minipage} \\
\midrule\noalign{}
\endhead
\bottomrule\noalign{}
\endlastfoot
Opening and recap & 10 min & Introductions, software check, workshop
map \\
Concept bridge & 25 min & R vs RStudio, why scripts and projects
matter \\
Guided coding & 60 min & Objects, vectors, indexing, \texttt{NA}, data
frames, summaries \\
Participant practice & 45 min & Apply the same workflow to a new vector
and data frame \\
Interpretation and reporting & 25 min & Turn a summary into a sentence a
colleague could read \\
Evidence log and capstone update & 15 min & Save the baseline task and
log it \\
\end{longtable}

\begin{center}\rule{0.5\linewidth}{0.5pt}\end{center}

\hypertarget{part-1-concept-bridge-core-workshop-activity}{%
\subsection{Part 1 -- Concept bridge (Core workshop
activity)}\label{part-1-concept-bridge-core-workshop-activity}}

\hypertarget{r-rstudio-scripts-and-the-console}{%
\subsubsection{R, RStudio, scripts and the
console}\label{r-rstudio-scripts-and-the-console}}

\begin{itemize}
\tightlist
\item
  \textbf{R} is the language and the engine that runs your code. It
  exists whether or not you ever open RStudio.
\item
  \textbf{RStudio} is an \emph{integrated development environment}
  (IDE): a window manager around R that adds a script editor, a file
  browser, a plot viewer and a help pane.
\item
  The \textbf{console} (bottom-left pane by default) runs one command at
  a time and forgets it once you close R. Use it for quick checks, not
  for anything you want to keep.
\item
  A \textbf{script} (a \texttt{.R} file) is a saved, repeatable sequence
  of commands. Everything that matters for your analysis should end up
  in a script, not only in the console.
\end{itemize}

\begin{quote}
\textbf{Working rule for this workshop:} if you would be unhappy to lose
it, it belongs in a script, not just in the console.
\end{quote}

\hypertarget{why-an-rstudio-project}{%
\subsubsection{Why an RStudio Project}\label{why-an-rstudio-project}}

An RStudio Project (\texttt{.Rproj} file) pins your working directory to
the project folder. This means:

\begin{itemize}
\tightlist
\item
  \texttt{read\_csv("Data/raw/file.csv")} works for you and for anyone
  else who opens the same project, on any computer.
\item
  You never need \texttt{setwd("C:/Users/yourname/Desktop/...")}, which
  breaks the moment the script moves to another machine.
\end{itemize}

\textbf{Today's setup steps:}

\begin{enumerate}
\def\labelenumi{\arabic{enumi}.}
\tightlist
\item
  \texttt{File\ -\textgreater{}\ New\ Project\ -\textgreater{}\ New\ Directory\ -\textgreater{}\ New\ Project}.
\item
  Name it something durable, e.g.~\texttt{r-workshop}.
\item
  Inside the project, create folders \texttt{Data/raw},
  \texttt{Data/processed}, \texttt{Scripts}, \texttt{Outputs} (you will
  reuse this structure through Day 7).
\item
  Create a new script \texttt{Scripts/day1\_baseline.R} and save it
  immediately.
\end{enumerate}

\begin{center}\rule{0.5\linewidth}{0.5pt}\end{center}

\hypertarget{part-2-guided-coding-core-workshop-activity}{%
\subsection{Part 2 -- Guided coding (Core workshop
activity)}\label{part-2-guided-coding-core-workshop-activity}}

Type each block into your script and run it line by line
(\texttt{Ctrl/Cmd\ +\ Enter}), not by pasting it all at once. Read the
output before moving on.

\hypertarget{assignment-object-names-and-comments}{%
\subsubsection{2.1 Assignment, object names and
comments}\label{assignment-object-names-and-comments}}

\begin{Shaded}
\begin{Highlighting}[]
\CommentTok{\# Comments start with \# and explain *why*, not *what the code obviously does*}
\NormalTok{workshop\_day }\OtherTok{\textless{}{-}} \DecValTok{1}          \CommentTok{\# use \textless{}{-} for assignment, not =}
\NormalTok{participant\_count }\OtherTok{\textless{}{-}} \DecValTok{24}
\NormalTok{average\_age }\OtherTok{\textless{}{-}} \FloatTok{41.6}

\NormalTok{workshop\_day}
\FunctionTok{print}\NormalTok{(participant\_count)}
\end{Highlighting}
\end{Shaded}

Object names should be readable in six months: \texttt{average\_age},
not \texttt{x}, not \texttt{avgAge1}.

\hypertarget{vectors}{%
\subsubsection{2.2 Vectors}\label{vectors}}

A vector is an ordered collection of values of the \emph{same} type.

\begin{Shaded}
\begin{Highlighting}[]
\NormalTok{ages }\OtherTok{\textless{}{-}} \FunctionTok{c}\NormalTok{(}\DecValTok{34}\NormalTok{, }\DecValTok{41}\NormalTok{, }\DecValTok{29}\NormalTok{, }\DecValTok{52}\NormalTok{, }\DecValTok{38}\NormalTok{, }\ConstantTok{NA}\NormalTok{, }\DecValTok{45}\NormalTok{)}
\FunctionTok{class}\NormalTok{(ages)}
\FunctionTok{length}\NormalTok{(ages)}

\NormalTok{regions }\OtherTok{\textless{}{-}} \FunctionTok{c}\NormalTok{(}\StringTok{"North West"}\NormalTok{, }\StringTok{"South East"}\NormalTok{, }\StringTok{"South West"}\NormalTok{, }\StringTok{"North Central"}\NormalTok{)}
\FunctionTok{class}\NormalTok{(regions)}
\end{Highlighting}
\end{Shaded}

\hypertarget{indexing-and-logical-conditions}{%
\subsubsection{2.3 Indexing and logical
conditions}\label{indexing-and-logical-conditions}}

\begin{Shaded}
\begin{Highlighting}[]
\NormalTok{ages[}\DecValTok{1}\NormalTok{]              }\CommentTok{\# first element}
\NormalTok{ages[}\FunctionTok{c}\NormalTok{(}\DecValTok{1}\NormalTok{, }\DecValTok{3}\NormalTok{)]        }\CommentTok{\# first and third}
\NormalTok{ages[ages }\SpecialCharTok{\textgreater{}} \DecValTok{40}\NormalTok{]      }\CommentTok{\# logical indexing: keep ages above 40}

\NormalTok{over\_40 }\OtherTok{\textless{}{-}}\NormalTok{ ages }\SpecialCharTok{\textgreater{}} \DecValTok{40} \CommentTok{\# this is itself a vector of TRUE/FALSE/NA}
\NormalTok{over\_40}
\end{Highlighting}
\end{Shaded}

Logical indexing is the single most useful R skill for the rest of this
workshop: almost every data-cleaning rule you write from Day 2 onward is
a logical condition like this one.

\hypertarget{missing-values}{%
\subsubsection{2.4 Missing values}\label{missing-values}}

\begin{Shaded}
\begin{Highlighting}[]
\FunctionTok{mean}\NormalTok{(ages)            }\CommentTok{\# returns NA {-}{-} R refuses to silently guess}
\FunctionTok{mean}\NormalTok{(ages, }\AttributeTok{na.rm =} \ConstantTok{TRUE}\NormalTok{)}
\FunctionTok{sum}\NormalTok{(}\FunctionTok{is.na}\NormalTok{(ages))      }\CommentTok{\# how many missing values?}
\FunctionTok{which}\NormalTok{(}\FunctionTok{is.na}\NormalTok{(ages))    }\CommentTok{\# which positions?}
\end{Highlighting}
\end{Shaded}

\texttt{NA} means ``unknown'', not zero and not ``missing on purpose''.
R propagates \texttt{NA} through almost every calculation until you tell
it otherwise with \texttt{na.rm\ =\ TRUE}. This is a safety feature, not
an inconvenience -- it stops you reporting a mean you did not actually
compute.

\hypertarget{a-small-data-frame}{%
\subsubsection{2.5 A small data frame}\label{a-small-data-frame}}

\begin{Shaded}
\begin{Highlighting}[]
\NormalTok{participants }\OtherTok{\textless{}{-}} \FunctionTok{data.frame}\NormalTok{(}
  \AttributeTok{participant\_id =} \FunctionTok{c}\NormalTok{(}\StringTok{"P01"}\NormalTok{, }\StringTok{"P02"}\NormalTok{, }\StringTok{"P03"}\NormalTok{, }\StringTok{"P04"}\NormalTok{, }\StringTok{"P05"}\NormalTok{),}
  \AttributeTok{age =} \FunctionTok{c}\NormalTok{(}\DecValTok{34}\NormalTok{, }\DecValTok{41}\NormalTok{, }\DecValTok{29}\NormalTok{, }\DecValTok{52}\NormalTok{, }\DecValTok{38}\NormalTok{),}
  \AttributeTok{region =} \FunctionTok{c}\NormalTok{(}\StringTok{"North West"}\NormalTok{, }\StringTok{"South East"}\NormalTok{, }\StringTok{"South West"}\NormalTok{, }\StringTok{"North Central"}\NormalTok{, }\StringTok{"South East"}\NormalTok{),}
  \AttributeTok{pre\_score =} \FunctionTok{c}\NormalTok{(}\DecValTok{58}\NormalTok{, }\DecValTok{64}\NormalTok{, }\DecValTok{71}\NormalTok{, }\DecValTok{49}\NormalTok{, }\DecValTok{60}\NormalTok{)}
\NormalTok{)}

\FunctionTok{str}\NormalTok{(participants)}
\FunctionTok{head}\NormalTok{(participants, }\DecValTok{3}\NormalTok{)}
\FunctionTok{nrow}\NormalTok{(participants)}
\FunctionTok{ncol}\NormalTok{(participants)}
\end{Highlighting}
\end{Shaded}

\hypertarget{descriptive-statistics-including-grouped-summaries}{%
\subsubsection{2.6 Descriptive statistics, including grouped
summaries}\label{descriptive-statistics-including-grouped-summaries}}

\begin{Shaded}
\begin{Highlighting}[]
\FunctionTok{mean}\NormalTok{(participants}\SpecialCharTok{$}\NormalTok{pre\_score)}
\FunctionTok{median}\NormalTok{(participants}\SpecialCharTok{$}\NormalTok{pre\_score)}
\FunctionTok{sd}\NormalTok{(participants}\SpecialCharTok{$}\NormalTok{pre\_score)}

\CommentTok{\# Grouped summary without extra packages}
\FunctionTok{tapply}\NormalTok{(participants}\SpecialCharTok{$}\NormalTok{pre\_score, participants}\SpecialCharTok{$}\NormalTok{region, mean)}
\end{Highlighting}
\end{Shaded}

We will redo grouped summaries with \texttt{dplyr::group\_by()} and
\texttt{summarise()} from Day 2 onward, because it scales better to many
groups and many statistics at once. \texttt{tapply()} is shown here so
the underlying idea -- split, compute, combine -- is clear before the
package hides the mechanics.

\hypertarget{reading-error-messages}{%
\subsubsection{2.7 Reading error
messages}\label{reading-error-messages}}

Run this deliberately:

\begin{Shaded}
\begin{Highlighting}[]
\FunctionTok{mean}\NormalTok{(participnts}\SpecialCharTok{$}\NormalTok{pre\_score)}
\end{Highlighting}
\end{Shaded}

R replies with something like
\texttt{Error\ in\ mean(participnts\$pre\_score)\ :\ object\ \textquotesingle{}participnts\textquotesingle{}\ not\ found}.
Read right to left: \emph{what} went wrong (object not found) and
\emph{where} (inside \texttt{mean(...)}). Most beginner errors are
typos, mismatched brackets/quotes, or using \texttt{=} where \texttt{==}
was meant. When stuck, run \texttt{?function\_name}
(e.g.~\texttt{?mean}) to open the help page, or
\texttt{example(function\_name)} for runnable examples.

\begin{center}\rule{0.5\linewidth}{0.5pt}\end{center}

\hypertarget{part-3-participant-practice-core-workshop-activity}{%
\subsection{Part 3 -- Participant practice (Core workshop
activity)}\label{part-3-participant-practice-core-workshop-activity}}

Work in your script. Do not look at the solution until you have tried
each step.

\begin{enumerate}
\def\labelenumi{\arabic{enumi}.}
\tightlist
\item
  Create a vector \texttt{attendance} with 8 values between 0 and 100,
  including one \texttt{NA}.
\item
  Compute the mean attendance, ignoring the \texttt{NA}.
\item
  Create a logical vector flagging attendance below 70.
\item
  Build a data frame with columns \texttt{student\_id} (made up),
  \texttt{attendance}, and \texttt{pass} (\texttt{TRUE}/\texttt{FALSE}
  for attendance \textgreater= 70, using your logical vector or an
  \texttt{ifelse()}).
\item
  Compute the proportion of students who pass.
\item
  Deliberately introduce one typo and read the resulting error message
  out loud to your neighbour before fixing it.
\end{enumerate}

A facilitator-reviewed solution is available in
\texttt{Solution-Scripts/day1\_practice\_solution.R} -- check it only
after attempting the task yourself.

\begin{center}\rule{0.5\linewidth}{0.5pt}\end{center}

\hypertarget{optional-demonstration}{%
\subsection{Optional demonstration}\label{optional-demonstration}}

These are useful if the group is moving quickly. They are not required
for the daily deliverable.

\begin{itemize}
\tightlist
\item
  \textbf{Matrices and basic matrix operations} -- \texttt{matrix()},
  \texttt{t()}, \texttt{\%*\%}. These return in full on Day 6 when you
  build a design matrix by hand.
\item
  \textbf{Factors} -- \texttt{factor()} for ordered or fixed categories,
  and why factors behave differently from plain character vectors in
  summaries and models.
\item
  \textbf{Lists and \texttt{typeof()}} -- a list can hold mixed types
  and different lengths; \texttt{typeof()} and \texttt{class()}
  distinguish the underlying storage from the object's behaviour.
\end{itemize}

\begin{Shaded}
\begin{Highlighting}[]
\CommentTok{\# Optional: factors}
\NormalTok{satisfaction }\OtherTok{\textless{}{-}} \FunctionTok{factor}\NormalTok{(}\FunctionTok{c}\NormalTok{(}\StringTok{"Low"}\NormalTok{, }\StringTok{"High"}\NormalTok{, }\StringTok{"Medium"}\NormalTok{, }\StringTok{"High"}\NormalTok{),}
                        \AttributeTok{levels =} \FunctionTok{c}\NormalTok{(}\StringTok{"Low"}\NormalTok{, }\StringTok{"Medium"}\NormalTok{, }\StringTok{"High"}\NormalTok{), }\AttributeTok{ordered =} \ConstantTok{TRUE}\NormalTok{)}
\NormalTok{satisfaction}
\NormalTok{satisfaction }\SpecialCharTok{\textless{}} \StringTok{"High"}
\end{Highlighting}
\end{Shaded}

\begin{center}\rule{0.5\linewidth}{0.5pt}\end{center}

\hypertarget{reference-or-take-home-material}{%
\subsection{Reference or take-home
material}\label{reference-or-take-home-material}}

\begin{itemize}
\tightlist
\item
  Extended debugging examples (mismatched parentheses, dangling commas,
  case-sensitive names).
\item
  More practice with subsetting: \texttt{{[}}, \texttt{{[}{[}},
  \texttt{\$}, and when each is required.
\item
  RStudio keyboard shortcuts: \texttt{Ctrl/Cmd\ +\ Enter} (run line),
  \texttt{Ctrl/Cmd\ +\ Shift\ +\ M} (pipe
  \texttt{\textbar{}\textgreater{}}), \texttt{Ctrl/Cmd\ +\ 1/2} (move
  between script and console).
\end{itemize}

\begin{center}\rule{0.5\linewidth}{0.5pt}\end{center}

\hypertarget{daily-deliverable-and-the-day-1-baseline-task}{%
\subsection{Daily deliverable and the Day 1 baseline
task}\label{daily-deliverable-and-the-day-1-baseline-task}}

\textbf{Deliverable:} a working RStudio project containing
\texttt{Scripts/day1\_baseline.R}, your practice exercise, and the
baseline task below, all saved inside the project folder.

\hypertarget{baseline-task-repeated-on-day-7}{%
\subsubsection{Baseline task (repeated on Day
7)}\label{baseline-task-repeated-on-day-7}}

A small dataset is provided at
\texttt{Data/raw/day1\_baseline\_sample.csv} (12 students, two columns
of interest plus an ID). It contains at least one missing value and at
least one invalid value -- you do not need to find every issue today,
just demonstrate the workflow.

\begin{enumerate}
\def\labelenumi{\arabic{enumi}.}
\tightlist
\item
  Import the dataset with \texttt{read.csv()} (we introduce
  \texttt{readr::read\_csv()} properly on Day 2).
\item
  Identify missing or clearly invalid values (for example, a negative
  score or an attendance percentage above 100) using the indexing skills
  from Part 2.
\item
  Produce one numerical summary (e.g.~mean test score, ignoring
  missing/invalid values) and one labelled plot (a basic \texttt{hist()}
  or \texttt{plot()} is enough today -- \texttt{ggplot2} starts Day 3).
\item
  Write three sentences interpreting the result for a non-technical
  reader (a department head, not a statistician).
\item
  Save the script, the plot, and your three sentences inside the project
  folder, e.g.~\texttt{Outputs/day1\_baseline\_summary.txt}.
\end{enumerate}

\begin{Shaded}
\begin{Highlighting}[]
\NormalTok{baseline }\OtherTok{\textless{}{-}} \FunctionTok{read.csv}\NormalTok{(}\StringTok{"Data/raw/day1\_baseline\_sample.csv"}\NormalTok{)}
\FunctionTok{str}\NormalTok{(baseline)}
\FunctionTok{summary}\NormalTok{(baseline}\SpecialCharTok{$}\NormalTok{test\_score)}
\FunctionTok{hist}\NormalTok{(baseline}\SpecialCharTok{$}\NormalTok{test\_score, }\AttributeTok{main =} \StringTok{"Day 1 baseline: test scores"}\NormalTok{, }\AttributeTok{xlab =} \StringTok{"Score"}\NormalTok{)}
\end{Highlighting}
\end{Shaded}

\hypertarget{evidence-log-entry}{%
\subsubsection{Evidence log entry}\label{evidence-log-entry}}

Open \texttt{Workshop-Evidence-Log.md} and complete the Day 1 row: note
the baseline task is saved, and that your R/RStudio setup is confirmed
working.

\begin{center}\rule{0.5\linewidth}{0.5pt}\end{center}

\hypertarget{facilitator-notes}{%
\subsection{Facilitator notes}\label{facilitator-notes}}

\begin{itemize}
\tightlist
\item
  Confirm every participant has R \textgreater= 4.x and a recent RStudio
  Desktop installed \emph{before} this session; software installation
  should not consume live-session time.
\item
  The Day 1 baseline task is intentionally small. Its only purpose is to
  give Day 7 something to compare against -- do not let it expand into a
  full analysis.
\item
  Keep ``Optional demonstration'' genuinely optional: if the group needs
  the full 60 minutes of guided coding plus 45 minutes of practice, skip
  the optional section entirely.
\end{itemize}
\end{document}
